\documentclass[11pt]{article}
\usepackage[T1]{fontenc}
\usepackage[utf8]{inputenc}
\usepackage[margin=1in]{geometry}
\usepackage{amsmath,amssymb,amsthm,microtype}
\usepackage[hidelinks]{hyperref}
\hypersetup{pdftitle={Equal dilatations deep in the Johnson filtration: a convex-core proof},pdfauthor={Alec Kriebel},pdfsubject={An attributed alternative proof of the negative answer to Farb Question 7.6; unrefereed}}
\newtheorem{theorem}{Theorem}
\newtheorem{lemma}[theorem]{Lemma}
\newtheorem{proposition}[theorem]{Proposition}
\theoremstyle{remark}\newtheorem*{remark}{Remark}
\newcommand{\Mod}{\operatorname{Mod}}
\newcommand{\Out}{\operatorname{Out}}
\newcommand{\Aut}{\operatorname{Aut}}
\newcommand{\tr}{\operatorname{tr}}
\newcommand{\PSL}{\operatorname{PSL}}
\newcommand{\SL}{\operatorname{SL}}
\newcommand{\HH}{\mathbb H^2}
\newcommand{\I}{\mathcal I}
\title{Equal dilatations deep in the Johnson filtration:\\a convex-core proof}
\author{Alec Kriebel\thanks{ORCID: \href{https://orcid.org/0009-0001-9320-500X}{0009-0001-9320-500X}.}}
\date{17 October 2026 \quad (version 1.0; unrefereed preprint)}
\begin{document}
\maketitle
\begin{abstract}
Farb asked whether the pseudo-Anosov dilatation spectrum of a fixed deep
Johnson-filtration term has bounded multiplicity, counted using ambient
mapping-class conjugacy. Ferudun has recorded a negative answer, including
this scope. We give an alternative proof using a compact convex core.
For every fixed closed genus $g\ge2$, fixed finite depth $k\ge3$, and $N$,
there are $N$ pairwise nonconjugate elements of $\I_g(k)$ in $\Mod_g$
with one common pseudo-Anosov dilatation. Explicit commutators produce
hyperbolic elements at the prescribed depth; high powers give a Schottky
subgroup. A compact-core argument uniformly bounds fusion of its conjugacy
classes in any discrete Fuchsian overgroup. Horowitz's universal trace
families then give the result. The argument does not require an arithmetic
description of the full disk stabilizer. No priority or novelty claim is made.
\end{abstract}

\section{Statement and context}
Let $S_g$ be a closed oriented surface of genus $g\ge2$, let
$\pi=\pi_1(S_g)$, and put
\[
 \Gamma_0=\pi,\qquad \Gamma_{j+1}=[\pi,\Gamma_j],\qquad
 \I_g(j)=\ker\bigl(\Mod_g\longrightarrow\Out(\pi/\Gamma_j)\bigr).
\]
Here $\Mod_g$ consists of orientation-preserving mapping classes.
Thus $\I_g(1)$ is the Torelli group. This is the indexing in
Farb's question \cite[\S1.2 and Question 7.6]{Farb}.
For $H\le\Mod_g$ and $t>0$, define the ambient multiplicity
\[
 m_H(t)=\#\{[f]_{\Mod_g}: f\in H\text{ is pseudo-Anosov and }
                         \log\lambda(f)=t\}.
\]
In this definition a class is counted when it has a representative in $H$.
Conjugators can belong to all of $\Mod_g$.

\begin{theorem}\label{thm:main}
For each fixed $g\ge2$ and fixed finite $k\ge3$, and every integer $N\ge1$,
there are pseudo-Anosov $f_1,\ldots,f_N\in\I_g(k)$ with the same dilatation
which represent distinct conjugacy classes of $\Mod_g$.
Consequently $\sup_{t>0}m_{\I_g(k)}(t)=\infty$.
\end{theorem}

This is the negative answer to Farb's Question 7.6. A theorem containing
this statement has already been published as an unrefereed preprint by
Ferudun \cite[Theorem 1.1]{Ferudun}; that theorem also states extensions
beyond the scope considered here. The present note supplies an alternative
proof of the displayed statement. Ferudun compares conjugacy using an
integral saddle-connection lattice and a virtually free group. Here we
first pass to a subgroup with compact convex core and use properness of
the full Fuchsian action. No assertion is made that this alternative
mechanism is new. The common dilatation may vary with $N$. Powers are
allowed; neither ambient primitivity nor the simple-length question
(Farb's Question 7.7) is addressed.

\section{A free subgroup at a fixed depth}
We begin by checking the filtration indices, including the outer action.
We use $[x,y]=xyx^{-1}y^{-1}$.

\begin{lemma}\label{lem:central}
For $a,b\ge1$ one has
$[\I_g(a),\I_g(b)]\subseteq\I_g(a+b)$.
A Dehn twist about a separating curve belongs to $\I_g(2)$.
\end{lemma}
\begin{proof}
Let $A(a)\le\Aut(\pi)$ consist of automorphisms inducing the identity on
$\pi/\Gamma_a$. The lower-central commutator inclusions give
$[\Gamma_i,\Gamma_j]\subseteq\Gamma_{i+j+1}$.
Induction on $j$ then gives
\begin{equation}\label{eq:action}
 \alpha(z)z^{-1}\in\Gamma_{a+j}
 \qquad(\alpha\in A(a),\ z\in\Gamma_j).
\end{equation}
Indeed, the case $j=0$ is the definition. At the next step expand
$\alpha([x,z])$ for $x\in\pi$, $z\in\Gamma_j$; the errors lie in
$[\Gamma_a,\Gamma_j]$ or $[\pi,\Gamma_{a+j}]$, hence in
$\Gamma_{a+j+1}$. Products and inverses preserve the congruence.

For $\alpha\in A(a)$, $\beta\in A(b)$ and $x\in\pi$, write
$\alpha(x)=xu$, $\beta(x)=xv$ with $u\in\Gamma_a$, $v\in\Gamma_b$.
By \eqref{eq:action}, modulo $\Gamma_{a+b}$ the values of
$\alpha\beta(x)$ and $\beta\alpha(x)$ are respectively $xuv$ and $xvu$.
They agree, since $[u,v]\in\Gamma_{a+b+1}$.
Thus $[A(a),A(b)]\subseteq A(a+b)$.
An outer automorphism trivial on $\pi/\Gamma_a$ has a representative
inducing an inner automorphism of that quotient. Lift its conjugating
element to $\pi$ and adjust the representative by the inverse inner
automorphism. This gives a representative in $A(a)$, and proves the
asserted containment after passage to outer classes. This centrality is
also the classical input used in \cite[proof of Proposition 4.8]{FLM}.

For a separating curve, take a basepoint on one side and a twist
representative which fixes that side. On generators on the other side
it acts by conjugation by the boundary word, up to its sign. That word
is in $\Gamma_1$, so its conjugation changes each generator by an element
of $[\Gamma_1,\pi]=\Gamma_2$. The induced outer action on
$\pi/\Gamma_2$ is therefore trivial.
\end{proof}

Fix filling separating simple closed curves $a,b$ in minimal position.
Such a pair exists for every $g\ge2$; see
\cite[\S7.2]{Farb} or \cite[Lemma 2.2]{AT}.
Write $\mu=i(a,b)$ and $A=T_a$, $B=T_b$.
The graph $a\cup b$ has $\mu$ vertices and $2\mu$ edges, and its
complement consists of $F\ge1$ disks. Thus
$\mu=2g-2+F\ge2g-1$. Since a separating curve has zero mod-$2$
intersection with every curve, $\mu$ is even and consequently $\mu\ge4$.

Thurston's rectangle construction gives a flat quadratic differential,
a Teichm\"uller disk $D$, and derivatives of $A,B$ with lifts
\begin{equation}\label{eq:AB}
 U=\begin{pmatrix}1&\mu\\0&1\end{pmatrix},\qquad
 V=\begin{pmatrix}1&0\\-\mu&1\end{pmatrix}.
\end{equation}
For a single pair of curves the shear parameter is the intersection
number $\mu$; see \cite[\S5.1 and Lemma 6.3]{Leininger}.
We use the following standard disk facts as stated in
\cite[Theorems 3.1--3.2]{Leininger}: the derivative homomorphism
\[
 \rho:\operatorname{Stab}_{\Mod_g}(D)\longrightarrow\PSL(2,\mathbb R)
\]
has discrete image and finite kernel; a hyperbolic derivative gives a
pseudo-Anosov with
\begin{equation}\label{eq:normalization}
 \log\lambda(f)=\tfrac12\ell_{\HH}(\rho(f));
\end{equation}
and a pseudo-Anosov stabilizes a unique Teichm\"uller disk.
The last statement is Bers's theorem, used here in Leininger's
restatement. We do not require injectivity of $\rho$ on the full stabilizer.

Define $c_1=A$ and $c_r=[B,c_{r-1}]$ for $r\ge2$.
Lemma~\ref{lem:central} gives $c_r\in\I_g(2r)$.
For any determinant-one matrix $C=\left(\begin{smallmatrix}p&q\\w&s\end{smallmatrix}\right)$,
direct multiplication gives
\[
 [V,C]=\begin{pmatrix}
 1+\mu qs&-\mu q^2\\
 \mu(s^2-1)-\mu^2qs&1-\mu qs+\mu^2q^2
 \end{pmatrix}.
\]
Its trace is $2+\mu^2q^2$. Starting with the upper-right entry $\mu$ of
$U$, induction yields, for the corresponding lifts $C_r$ of $\rho(c_r)$,
\begin{equation}\label{eq:deep}
 (C_r)_{12}=-\mu^{\,2^r-1},\qquad
 \tr C_r=2+\mu^{\,2^r}>2 \qquad(r\ge2).
\end{equation}
Fix $k\ge3$ and set $r=\max\{2,\lceil k/2\rceil\}$ and $c=c_r$.
Then $c\in\I_g(k)$ and its derivative is hyperbolic.
Its two fixed points on $\partial\HH=\mathbb R\cup\{\infty\}$ avoid $0$,
since $(C_r)_{12}\ne0$. The parabolic
$V^M(z)=z/(1-M\mu z)$ sends both points towards $0$ as $M\to\infty$.
For some positive $M$, the fixed points of $\rho(c)$ and
$\rho(d)$, where $d=B^McB^{-M}$, are therefore disjoint.
Normality gives $d\in\I_g(k)$.

\begin{lemma}\label{lem:schottky}
Two hyperbolic elements of $\PSL(2,\mathbb R)$ with disjoint fixed-point
sets have sufficiently large common positive powers which freely
generate a convex-cocompact rank-two group.
\end{lemma}
\begin{proof}
Choose disjoint closed boundary intervals around their four fixed
points. Choose a repelling geodesic half-plane for each element inside
its repelling interval. For sufficiently large powers, the image of
the complementary half-plane is an attracting half-plane inside the
corresponding attracting interval. The four half-planes have disjoint
closures, including their ideal endpoints. Their pairings give the
classical Schottky fundamental region $P$, exterior to their interiors,
and ping-pong proves that the two powered elements are a free basis.

Here strict separation also gives compactness of the convex core.
For each possible first letter, the limit set lies in the images under
that letter of the three allowed next-letter intervals. These finitely
many images are strictly inside the corresponding paired interval.
Thus the ideal boundary of the closed fundamental region $P$ is disjoint
from the limit set $\Lambda$. The closed convex hull $X$ of $\Lambda$
in $\HH$ has ideal boundary exactly $\Lambda$. Consequently $P\cap X$
has no ideal accumulation point and is compact. Its translates cover
$X$, so $X/J$ is compact for the generated group $J$. In particular all
its nonidentity elements are hyperbolic.
\end{proof}

Apply the lemma to $\rho(c),\rho(d)$ and choose the resulting exponent
$L$. Then
\begin{equation}\label{eq:H}
 H=\langle c^L,d^L\rangle\le\I_g(k),\qquad J=\rho(H)
\end{equation}
is free of rank two and $J$ is convex cocompact.
Indeed, any relation between the displayed generators of $H$ would
produce one between their free images; this also proves that
$\rho|_H:H\to J$ is an isomorphism. Every nonidentity element of $H$
is pseudo-Anosov. The choices of $H,J$ are now fixed independently of $N$.

\section{A uniform bound for ambient conjugacy}
The following elementary lemma is the conjugacy comparison needed here.
Convex cocompact means that the quotient of the convex hull of the limit
set is compact.

\begin{lemma}\label{lem:fusion}
Let $J\le G\le\PSL(2,\mathbb R)$, where $G$ is discrete and $J$ is
convex cocompact. There is a finite constant $C=C(J,G)$ such that each
$G$-conjugacy class of hyperbolic elements contains at most $C$ distinct
$J$-conjugacy classes of elements of $J$.
\end{lemma}
\begin{proof}
Let $X$ be the convex hull of $\Lambda(J)$, and choose a compact set
$K\subset X$ with $JK=X$. Such a set exists by compactness of $X/J$.
Properness of the discrete action on $\HH$ implies finiteness of
\[
 S=\{s\in G:sK\cap K\ne\varnothing\}.
\]
Explicitly, if $K\subset B(o,R)$ then $d(o,so)\le2R$ for $s\in S$;
the isometries satisfying this bound form a compact set because the
stabilizer of $o$ is compact. Its intersection with the discrete group
$G$ is finite.

Fix a hyperbolic $u\in J$. Its axis is in $X$. Conjugate $u$ within $J$
so that a point $x$ on its axis belongs to $K$. Keep this one $u$ and this
one $x$ fixed for all subsequent conjugators. If $v=tut^{-1}\in J$ with
$t\in G$, then $tx$ is on the axis of $v$, hence in $X$.
Choose $j\in J$ with $jtx\in K$. Now $s=jt$ belongs to $S$ and
\[
 jvj^{-1}=sus^{-1}.
\]
Thus every $J$-class in the specified $G$-class is represented by an
entry of the fixed list $\{sus^{-1}:s\in S\}$ that belongs to $J$.
Take $C=|S|$. Although $x$ depends on $u$, the sets $K,S$ do not;
the bound is uniform in $u$ and its translation length.
\end{proof}

Let $G=\rho(\operatorname{Stab}_{\Mod_g}(D))$ and take $J$ from
\eqref{eq:H}. If $f,h\in H\setminus\{1\}$ satisfy $zfz^{-1}=h$
for any $z\in\Mod_g$, then $h$ stabilizes both $D$ and $zD$.
Its unique invariant disk forces $zD=D$. Consequently
$\rho(f),\rho(h)$ are conjugate in the full group $G$.
Since $H\to J$ is an isomorphism, Lemma~\ref{lem:fusion} bounds by
$C$ the number of $H$-classes inside one ambient $\Mod_g$-class.
The possibly nontrivial finite kernel of the full derivative map causes
no loss: only injectivity on $H$ and the implication from ambient
conjugacy to $G$-conjugacy were used. Neither finite index nor arithmetic
of $G$ is required.

\section{Universal trace families and the conclusion}
Horowitz \cite[Example 8.2, pp.~648--649]{Horowitz} constructs, in a free
group $F(x,y)$, $2^m$ nonidentity pairwise nonconjugate words with the
same trace under every representation to $\SL(2,\mathbb C)$.
For clarity, his recursion can be written as
\begin{align}\label{eq:horowitz}
 W_0&=x,\notag\\
 W_j&=W_{j-1}^{-\varepsilon_j}y^{2j}
       W_{j-1}^{\varepsilon_j}y^{2j-1}
       W_{j-1}^{-\varepsilon_j}y^{2j}
       W_{j-1}^{\varepsilon_j},\quad \varepsilon_j\in\{1,-1\}.
\end{align}
The exponents $2j$ are linear in $j$; the family size is $2^m$.
We recall why both properties of this construction hold.

The words are freely reduced and begin and end with $x^{\pm1}$.
For $U=W_{m-1}$ the cyclic form is
\[
 U^{\varepsilon_m}y^{2m-1}U^{-\varepsilon_m}y^{4m}.
\]
The positive odd syllable $y^{2m-1}$ occurs exactly once, and its inverse
syllable is absent; earlier words have only smaller odd syllables.
The closing syllable $y^{4m}$ is also unique, since every syllable
in $U$ has absolute exponent at most $2m-2$. These two markers
anchor cyclic comparison and recover the block $U^{\varepsilon_m}$. For $m>1$ the corresponding odd marker
in $U$ distinguishes $U$ from the inverse of any earlier-level choice;
for $m=1$ use $x\ne x^{-1}$. Induction recovers all the signs.
Thus the $2^m$ choices give distinct conjugacy classes, and the words
are nonidentity.

For the trace assertion, the rank-two trace-polynomial identity says
that the trace of a word in $P,Q\in\SL(2,\mathbb C)$ is a polynomial
in $\tr P,\tr Q,\tr(PQ)$
\cite[Theorem 3.1 and Corollary 3.14]{Horowitz}.
Flipping the sign at the $j$th stage replaces the pair
$(U^{-1}yU,y)$ by $(UyU^{-1},y)$, where $U=W_{j-1}$; every later
word is the same word in the respective pair. Their trace coordinates
agree, since
\[
 \tr(U^{-1}yU)=\tr(UyU^{-1})=\tr y,\qquad
 \tr(U^{-1}yUy)=\tr(UyU^{-1}y).
\]
The latter equality is cyclic invariance of trace. Flipping one sign
at a time proves the universal common-trace assertion.

Identify $F(x,y)$ with $H$ using the basis in \eqref{eq:H}.
The projective representation of $H$ lifts to $\SL(2,\mathbb R)$ by
choosing lifts of the two free generators. At level $m$, the $2^m$
words therefore have a common trace $t_m$. They are all hyperbolic,
so $|t_m|>2$. Their common translation length and dilatation satisfy
\[
 \ell_m=2\operatorname{arcosh}(|t_m|/2),\qquad
 \lambda_m=\exp(\ell_m/2)
           =\frac{|t_m|+\sqrt{t_m^2-4}}2>1
\]
by \eqref{eq:normalization}. At most $C$ of these $H$-classes can fuse
in a $\Mod_g$-class. There are consequently at least
$\lceil2^m/C\rceil$ ambient classes at the common value
$\log\lambda_m$. Given $N$, choose $m$ with $2^m\ge CN$.
This proves Theorem~\ref{thm:main}, with the genus, depth and subgroup
all fixed before $m$ varies. \hfill$\square$

\section*{Scope, verification and disclosure}
The argument includes genus two and each finite Johnson depth
separately. It asserts no infinite multiplicity at one fixed value,
no result for the intersection of all filtration terms, and no
ambient-primitivity or simple-loop statement. Its infinite claims rest
on the written arguments and the cited standard theorems. Accompanying
exact-arithmetic checks of commutator matrices and finite trace-word
families are consistency checks, not substitutes for those proofs.

AI tools were used extensively in research, proof development,
verification, literature searches and writing. Additional AI agents
performed adversarial mathematical checks. This manuscript is
unrefereed, has not undergone external human peer review, and is not
a formal proof-assistant certification. The author is responsible for
the claims and final text. No claim of first resolution, a new existence
theorem, or current openness of the historical question is made.

\begin{thebibliography}{9}
\bibitem{AT} T.~Aougab and S.~J.~Taylor,
\emph{Pseudo-Anosovs optimizing the ratio of Teichm\"uller to curve graph
translation length}, arXiv:1510.00995, Lemma~2.2, preprint pp.~3--4.
\url{https://arxiv.org/abs/1510.00995}.

\bibitem{Farb} B.~Farb,
\emph{Some problems on mapping class groups and moduli space},
in \emph{Problems on Mapping Class Groups and Related Topics},
Proc. Sympos. Pure Math. \textbf{74}, American Mathematical Society,
2006, pp.~11--55.
\href{https://doi.org/10.1090/pspum/074/2264130}{doi:10.1090/pspum/074/2264130}.
Question~7.6 and the filtration conventions are cited here in the
author-hosted volume version, printed pp.~51 and~14 respectively
(PDF pp.~58 and~21):
\url{https://www.math.uchicago.edu/~farb/papers/mcgbook.pdf}.
Its pagination differs from the published chapter.

\bibitem{FLM} B.~Farb, C.~J.~Leininger and D.~Margalit,
\emph{The lower central series and pseudo-Anosov dilatations},
Amer. J. Math. \textbf{130} (2008), 799--827.
\href{https://doi.org/10.1353/ajm.0.0005}{doi:10.1353/ajm.0.0005}.
Preprint: \url{https://arxiv.org/abs/math/0603675},
\S4.2, proof of Proposition~4.8.

\bibitem{Ferudun} A.~Ferudun,
\emph{Unbounded Multiplicity in the Dilatation Spectrum of Every Term
of the Johnson Filtration: An Observation on a Question of Farb},
EulerSolve Research Papers, AMR-109-0082, version~1.0,
11 October 2026, unrefereed preprint.
\href{https://doi.org/10.5281/zenodo.23301731}{doi:10.5281/zenodo.23301731}.

\bibitem{Horowitz} R.~D.~Horowitz,
\emph{Characters of free groups represented in the two-dimensional
special linear group}, Comm. Pure Appl. Math. \textbf{25} (1972),
635--649.
\href{https://doi.org/10.1002/cpa.3160250602}{doi:10.1002/cpa.3160250602}.
Primary paper copy:
\url{https://www.math.purdue.edu/~eremenko/Pdf/horowitz.pdf}.

\bibitem{Leininger} C.~J.~Leininger,
\emph{On groups generated by two positive multi-twists:
Teichm\"uller curves and Lehmer's number},
Geom. Topol. \textbf{8} (2004), 1301--1359.
\href{https://doi.org/10.2140/gt.2004.8.1301}{doi:10.2140/gt.2004.8.1301}.
Theorems~3.1--3.2, pp.~1314--1315; \S5.1 and Lemma~6.3.
\end{thebibliography}
\end{document}
